\documentclass[11pt,a4paper]{article}
\usepackage[utf8]{inputenc}
\usepackage[T1]{fontenc}
\usepackage[portuguese]{babel}
\usepackage{geometry}
\geometry{margin=2.5cm}
\usepackage{listings}
\usepackage{xcolor}
\usepackage{amsmath}
\usepackage{amssymb}
\usepackage{hyperref}
\usepackage{enumitem}
\usepackage{tcolorbox}
\usepackage{tabularx}
\usepackage{booktabs}
\usepackage{graphicx}

\definecolor{codebg}{rgb}{0.95,0.95,0.95}
\definecolor{codegreen}{rgb}{0.1,0.5,0.1}
\definecolor{codepurple}{rgb}{0.5,0.1,0.5}
\definecolor{codegray}{rgb}{0.5,0.5,0.5}

\lstdefinestyle{python}{
    backgroundcolor=\color{codebg},
    commentstyle=\color{codegray}\itshape,
    keywordstyle=\color{codepurple}\bfseries,
    stringstyle=\color{codegreen},
    basicstyle=\ttfamily\small,
    breaklines=true,
    frame=single,
    language=Python,
    showstringspaces=false,
    tabsize=4,
    literate={á}{{\'a}}1 {ã}{{\~a}}1 {é}{{\'e}}1 {í}{{\'i}}1 {ó}{{\'o}}1 {õ}{{\~o}}1 {ú}{{\'u}}1 {ç}{{\c{c}}}1 {Á}{{\'A}}1 {Ã}{{\~A}}1 {É}{{\'E}}1 {Í}{{\'I}}1 {Ó}{{\'O}}1 {Õ}{{\~O}}1 {Ú}{{\'U}}1 {Ç}{{\c{C}}}1 {â}{{\^a}}1 {ê}{{\^e}}1 {ô}{{\^o}}1 {û}{{\^u}}1 {Â}{{\^A}}1 {Ê}{{\^E}}1 {Ô}{{\^O}}1 {Û}{{\^U}}1 {à}{{\`a}}1 {è}{{\`e}}1 {ì}{{\`i}}1 {ò}{{\`o}}1 {ù}{{\`u}}1 {ä}{{\"a}}1 {ö}{{\"o}}1 {Ä}{{\"A}}1 {Ö}{{\"O}}1
}
\lstset{style=python}

\newtcolorbox{objetivos}[1][]{
  colback=green!5!white,
  colframe=green!40!black,
  title=O que você vai aprender nesta página,
  fonttitle=\bfseries,
  #1
}

\newtcolorbox{exercicio}[1]{
  colback=blue!5!white,
  colframe=blue!40!black,
  title=#1,
  fonttitle=\bfseries
}

\newtcolorbox{solucao}{
  colback=yellow!5!white,
  colframe=orange!60!black,
  title=Solução proposta,
  fonttitle=\bfseries
}

\title{\textbf{Data Analysis with Python 2024--2025}\\Parte 6: Machine Learning Types\\\large Soluções dos Exercícios}
\author{Tradução e Adaptação: Universidade de Helsinque (MOOC.fi)}
\date{}

\begin{document}

\maketitle

\section*{Nota Importante}
Este material é uma tradução e adaptação baseada no curso \textit{Data Analysis with Python 2024-2025}, oferecido pela \textbf{Universidade de Helsinque (MOOC.fi)}. As soluções apresentadas a seguir foram elaboradas para fins didáticos, seguindo os enunciados dos exercícios propostos no curso original.

\tableofcontents
\newpage

\section{Soluções - Capítulo 6 (Clustering)}

\subsection*{Exercício 6.5 -- Clustering de Plantas}

\begin{solucao}
\begin{lstlisting}
import scipy
from sklearn.datasets import load_iris
from sklearn.cluster import KMeans
from sklearn.metrics import accuracy_score

def find_permutation(n_clusters, real_labels, labels):
    permutation = []
    for i in range(n_clusters):
        idx = labels == i
        # Escolhe o rotulo mais frequente mapeado no escopo do cluster
        new_label = scipy.stats.mode(real_labels[idx])[0]
        # Validacao cross-version pandas/scipy
        if type(new_label) == np.ndarray:
            new_label = new_label[0]
        permutation.append(new_label)
    return permutation

def plant_clustering():
    iris = load_iris()
    X = iris.data
    y = iris.target
    
    # Modelo preve em 3 categorias independentes
    model = KMeans(n_clusters=3, random_state=0)
    model.fit(X)
    
    # Realinha a resposta com a chave e submete ao validador
    permutation = find_permutation(3, y, model.labels_)
    new_labels = [permutation[label] for label in model.labels_]
    
    return accuracy_score(y, new_labels)
\end{lstlisting}
\end{solucao}

\subsection*{Exercício 6.6 -- Clusters Não Convexos}

\begin{solucao}
\begin{lstlisting}
import pandas as pd
import numpy as np
import scipy
from sklearn.cluster import DBSCAN
from sklearn.metrics import accuracy_score

def find_permutation_with_outliers(n_clusters, real_labels, labels):
    permutation = []
    for i in range(n_clusters):
        idx = labels == i
        if len(real_labels[idx]) == 0:
            permutation.append(0) # Padrao caso nao haja registros
            continue
        new_label = scipy.stats.mode(real_labels[idx])[0]
        if type(new_label) == np.ndarray:
            new_label = new_label[0]
        permutation.append(new_label)
    return permutation

def nonconvex_clusters():
    df = pd.read_csv("src/data.tsv", sep='\t')
    X = df[['X1', 'X2']].values
    y = df['y'].values
    
    results = []
    for eps in np.arange(0.05, 0.20, 0.05):
        model = DBSCAN(eps=eps)
        model.fit(X)
        
        # Filtra quantidade e outliers (-1)
        unique_labels = set(model.labels_)
        num_clusters = len(unique_labels) - (1 if -1 in unique_labels else 0)
        outliers = list(model.labels_).count(-1)
        
        # Verificando se os espacos convergem com os reais previstos (2 clusters isolados)
        if num_clusters == len(set(y)):
            perm = find_permutation_with_outliers(num_clusters, y, model.labels_)
            mask = model.labels_ != -1
            acc = accuracy_score(y[mask], [perm[label] for label in model.labels_[mask]])
        else:
            acc = np.nan
            
        results.append([eps, acc, num_clusters, outliers])
        
    return pd.DataFrame(results, columns=["eps", "Score", "Clusters", "Outliers"])
\end{lstlisting}
\end{solucao}

\subsection*{Exercício 6.7 -- Locais de Ligação Molecular}

\begin{solucao}
\begin{lstlisting}
import pandas as pd
import numpy as np
import scipy
from sklearn.cluster import AgglomerativeClustering
from sklearn.metrics import accuracy_score, pairwise_distances

def toint(nucleotide):
    mapping = {'A': 0, 'C': 1, 'G': 2, 'T': 3}
    return mapping.get(nucleotide, -1)

def get_features_and_labels(filename):
    df = pd.read_csv(filename, sep='\t')
    
    # List comprehension encadeada gerando vetores inteiros
    X = np.array([[toint(char) for char in seq] for seq in df['X']])
    y = df['y'].values
    return X, y

def find_permutation(n_clusters, real_labels, labels):
    permutation = []
    for i in range(n_clusters):
        idx = labels == i
        new_label = scipy.stats.mode(real_labels[idx])[0]
        if type(new_label) == np.ndarray:
            new_label = new_label[0]
        permutation.append(new_label)
    return permutation

def cluster_euclidean(filename):
    X, y = get_features_and_labels(filename)
    # Classificador aglomerativo Euclidiano
    model = AgglomerativeClustering(n_clusters=2, linkage='average', affinity='euclidean')
    labels = model.fit_predict(X)
    
    perm = find_permutation(2, y, labels)
    new_labels = [perm[l] for l in labels]
    return accuracy_score(y, new_labels)

def cluster_hamming(filename):
    X, y = get_features_and_labels(filename)
    
    # Processa os componentes metricos usando algorimos focados no modelo Hamming
    dist_matrix = pairwise_distances(X, metric='hamming')
    
    # Utiliza as distancias calculadas previamente via precomputed
    model = AgglomerativeClustering(n_clusters=2, linkage='average', affinity='precomputed')
    labels = model.fit_predict(dist_matrix)
    
    perm = find_permutation(2, y, labels)
    new_labels = [perm[l] for l in labels]
    return accuracy_score(y, new_labels)
\end{lstlisting}
\end{solucao}

\end{document}
